\section{Verificación Processor-in-the-Loop}

En PIL, el código del controlador se compila para el procesador objetivo y se ejecuta en él,
mientras la planta sigue simulada en la PC \cite{mw_pil}. El procesador es una ESP32 (ESP32-D0WD-V3,
placa de desarrollo de 38 pines), la misma que ejecuta el controlador en el prototipo.

El \emph{Simulink Support Package for Arduino Hardware} de R2026a programa la ESP32, pero no
habilita PIL para ella \cite{mw_arduino}: la clase de la placa declara
\texttt{IsPilSupported = 0}, y al registrar a mano una conexión PIL, la clase interna del paquete
falla porque espera datos que solo existen para las placas soportadas. Por eso se usó una
alternativa de ingeniería: un PIL propio con la misma arquitectura que el de Simulink. En la ESP32 corre exactamente el código C de la Etapa~E, el mismo
verificado en SIL (se comprueba que la copia sea idéntica byte a byte), compilado con la
cadena Xtensa del paquete, con dos instancias, una por eje. Un sketch hace solo de puente serie:
recibe $r$, $\xh$, $\vh$ y \texttt{habilitado} como \texttt{double} de \SI{8}{bytes}, llama a
\texttt{Controller\_BoP\_step}, y devuelve $\thc$, $I$, \texttt{saturado} y los ciclos de CPU de
la llamada. En la PC siguen la planta del banco SIL, la referencia, la rama Normal y la
comparación.

La Fig.~\ref{fig:pilarq} resume la arquitectura. El intercambio es por el puerto USB a
\SI{115200}{baudios}, un paso por mensaje, sin restricción de tiempo real.

\begin{figure}[!t]
\centering
\begin{tikzpicture}[font=\scriptsize, >=stealth,
  caja/.style={draw, rounded corners=1pt, align=center, minimum height=0.75cm, inner sep=2pt},
  grupo/.style={draw, dashed, rounded corners=2pt, inner sep=5pt}]
  \node[caja] (ref) {Referencia\\y habilitación};
  \node[caja, below=0.45cm of ref] (planta) {Planta de la rama PIL\\(2 ejes, Kalman)};
  \node[caja, below=0.45cm of planta] (normal) {Rama Normal\\(modelo y su planta)};
  \node[caja, right=0.4cm of normal] (cmp) {Comparación\\muestra a muestra};
  \node[grupo, fit=(ref)(planta)(normal)(cmp), label={[font=\scriptsize]above:PC (MATLAB)}] (pc) {};
  \node[caja, right=2.7cm of planta, minimum width=2.1cm] (step)
       {Puente serie\\$\downarrow$\\\texttt{Controller\_BoP\_step}\\(2 instancias)};
  \node[grupo, fit=(step), label={[font=\scriptsize]above:ESP32}] (esp) {};
  \draw[->] ([yshift=4pt]planta.east) -- node[above, font=\tiny, yshift=1pt]{$r, \xh, \vh$, hab.} ([yshift=4pt]step.west);
  \draw[<-] ([yshift=-4pt]planta.east) -- node[below, font=\tiny, yshift=-1pt]{$\thc, I$, sat.} ([yshift=-4pt]step.west);
  \draw[->] (ref) -- (planta);
  \draw[->] (normal) -- (cmp);
  \draw[->] (planta.south east) -- (cmp.north);
  \draw[->] (ref.west) -- ([xshift=-3mm]ref.west -| planta.west) |- (normal.west);
\end{tikzpicture}
\caption{Arquitectura del PIL propio. En la ESP32 corre solo el código generado; la planta, la
referencia, la rama Normal y la comparación siguen en la PC.}
\label{fig:pilarq}
\end{figure}

Se hicieron dos pruebas. En G4 se ejecutó el controlador solo, con las entradas del MIL de los
casos D1 y D2; el segundo tiene nueve cuadros saturados, así que ejercita la saturación y el
\emph{clamping} que el ensayo SIL no alcanzaba. En los dos casos $\thc$ e $I$ fueron idénticos y
ninguna bandera de saturación difirió. En G5 se repitió el banco de lazo cerrado de la
Sección~IV-C con la rama PIL. La Fig.~\ref{fig:pil} muestra las diferencias y la
Tabla~\ref{tab:silpil} resume las dos verificaciones.

\begin{figure}[!t]
\centering
\begin{tikzpicture}
\begin{axis}[estilo, height=3.2cm, xmin=0, xmax=9, ymin=-1, ymax=1,
             xlabel={Tiempo (\si{\second})}, ylabel={Diferencia}, legend pos=north east,
             legend columns=2]
  \addplot[cRojo] table[col sep=comma, x=t_s, y=dif_thx] {\G/G5_comparacion.csv};
  \addplot[cVerde, dashed] table[col sep=comma, x=t_s, y=dif_x] {\G/G5_comparacion.csv};
  \legend{$\Delta\theta_x$ (\si{\radian}), $\Delta x$ (\si{\milli\metre})}
\end{axis}
\end{tikzpicture}
\caption{Diferencias PIL $-$ Normal en el banco de lazo cerrado (G5).}
\label{fig:pil}
\end{figure}

\begin{table}[!t]
\caption{Resultados de las verificaciones SIL y PIL. LC: banco de lazo cerrado; G4: controlador
solo con las entradas del MIL ($\thc$, $I$ y \texttt{saturado}). Fuente:
\texttt{etapaF/resultados/F1\_metricas.csv}, \texttt{etapaG/resultados/G4\_resultados.csv} y
\texttt{G5\_metricas.csv}.}
\label{tab:silpil}
\centering
\footnotesize
\setlength{\tabcolsep}{3pt}
\begin{tabular}{@{}lccc@{}}
\toprule
 & \textbf{SIL (F1)} & \textbf{PIL (G4)} & \textbf{PIL (G5)} \\
\midrule
Ensayo & LC, 2 ejes & D1 y D2 & LC, 2 ejes \\
Muestras & 271 & 194 y 179 & 271 \\
Señales & 10 & 3 & 10 \\
$\max|\Delta|$ & 0 & 0 & 0 \\
RMSE & 0 & 0 & 0 \\
Igualdad exacta & 271/271 & Sí & 271/271 \\
Cuadros saturados & 0 & 9 (en D2) & 0 \\
\midrule
Resultado & PASS & PASS & PASS \\
\bottomrule
\end{tabular}
\end{table}

La igualdad exacta indica que el compilador Xtensa y la aritmética \texttt{double} de la ESP32,
que no tiene unidad de punto flotante de doble precisión y la emula por software, dan los mismos
resultados bit a bit que la PC para esta secuencia de operaciones. El tiempo de
\texttt{Controller\_BoP\_step} en la ESP32 fue de \SI{7.5}{\micro\second} de mediana y
\SI{31.3}{\micro\second} como máximo, el \SI{0.094}{\percent} de $\Ts$, así que el cálculo no
limita el periodo de muestreo.

PIL aporta, sobre SIL, la evidencia de que el código corre en el procesador objetivo, con su
compilador, y conserva el comportamiento. Aun así, un PASS en PIL no valida el controlador sobre
la planta real: la planta sigue siendo un modelo, no hay cámara ni servos, el intercambio por el
puerto serie no corre a \SI{30}{fps} en tiempo real, y el modelo usado para diseñar resultó, como
se vio en la Sección~II, bastante distinto del prototipo.
